% =========================================================
% Computational Linear Algebra — Rules for Matrix Operations
% (Strang §2.4 / Johnston §1.3)
% =========================================================

\section{Rules for Matrix Operations}

\begin{tcolorbox}[colback=gray!6, colframe=gray!40, arc=2pt,
  left=6pt, right=6pt, top=4pt, bottom=4pt,
  title={\small\textbf{Matrix Operations --- Quick Reference}},
  fonttitle=\small\bfseries]
\begin{tabular}{@{}p{0.30\textwidth}p{0.63\textwidth}@{}}
\textbf{Associativity}    & $A(BC) = (AB)C$. \\[4pt]
\textbf{Distributivity}   & $A(B+C) = AB + AC$; $(A+B)C = AC + BC$. \\[4pt]
\textbf{NOT commutative}  & $AB \neq BA$ in general. \\[4pt]
\textbf{Identity}         & $AI = IA = A$. \\[4pt]
\textbf{Transpose product}& $(AB)^T = B^T A^T$. \\[4pt]
\textbf{Diagonal matrices}& $D_1 D_2$ is diagonal; $DA$ scales rows; $AD$ scales columns. \\
\end{tabular}
\end{tcolorbox}

% ---------------------------------------------------------
\subsection{Algebraic Rules}

\begin{tcolorbox}[colback=thmbox, colframe=thmborder, arc=2pt,
  left=6pt, right=6pt, top=4pt, bottom=4pt,
  title={\small\textbf{Theorem (Rules for Matrix Multiplication)}},
  fonttitle=\small\bfseries]
When the dimensions are compatible:
\begin{enumerate}
  \item $A(BC) = (AB)C$ \hfill (associativity)
  \item $A(B + C) = AB + AC$ \hfill (left distributivity)
  \item $(A + B)C = AC + BC$ \hfill (right distributivity)
  \item $c(AB) = (cA)B = A(cB)$ for any scalar $c$
  \item $(AB)^T = B^T A^T$ \hfill (transpose reverses order)
\end{enumerate}
The following do \textbf{not} hold in general:
\begin{itemize}
  \item $AB \neq BA$ (not commutative)
  \item $AB = AC \not\Rightarrow B = C$ (no cancellation in general)
  \item $AB = 0 \not\Rightarrow A = 0$ or $B = 0$
\end{itemize}
\end{tcolorbox}

% ---------------------------------------------------------
\subsection{Diagonal Matrices}

\begin{tcolorbox}[colback=propbox, colframe=propborder, arc=2pt,
  left=6pt, right=6pt, top=4pt, bottom=4pt,
  title={\small\textbf{Definition (Diagonal Matrix)}},
  fonttitle=\small\bfseries]
A \emph{diagonal matrix} $D \in \mathbb{R}^{n \times n}$ has $D_{ij} = 0$
for all $i \neq j$.
Written $D = \operatorname{diag}(d_1, d_2, \ldots, d_n)$.

Key computational facts:
\begin{itemize}
  \item $DA$ multiplies each \emph{row} $i$ of $A$ by $d_i$.
  \item $AD$ multiplies each \emph{column} $j$ of $A$ by $d_j$.
  \item $D_1 D_2 = \operatorname{diag}(d_1 e_1, \ldots, d_n e_n)$ where $e_j$ are the diagonal entries of $D_2$.
\end{itemize}
\end{tcolorbox}

% ---------------------------------------------------------
\subsection{Block Matrices}

\begin{remark*}[Block matrix multiplication]
Matrices can be partitioned into blocks, and multiplied block-by-block
as if the blocks were scalar entries, provided the block sizes are
compatible.
If $A = \begin{pmatrix} A_{11} & A_{12} \\ A_{21} & A_{22} \end{pmatrix}$
and $B = \begin{pmatrix} B_{11} & B_{12} \\ B_{21} & B_{22} \end{pmatrix}$,
then
\[
  AB = \begin{pmatrix}
    A_{11}B_{11} + A_{12}B_{21} & A_{11}B_{12} + A_{12}B_{22} \\
    A_{21}B_{11} + A_{22}B_{21} & A_{21}B_{12} + A_{22}B_{22}
  \end{pmatrix},
\]
provided inner block dimensions match.
The standard matrix of a linear transformation $[T] = [T(\mathbf{e}_1) \mid T(\mathbf{e}_2) \mid \cdots \mid T(\mathbf{e}_n)]$
is a block matrix with column vectors as blocks, and
$[T]\mathbf{v} = v_1 T(\mathbf{e}_1) + \cdots + v_n T(\mathbf{e}_n)$
is block matrix--vector multiplication.
\end{remark*}
